This book is concerned with the concept trace formulas with a focus on trace formulas of automorphic forms on $\textrm{SL}_2(\mathbb{Z})$. In general a trace formula is a formula where the trace of some operator is computed in two different ways leading to an equality. We call these the spectral side and geometric side of the trace formula. 

This chapter establishes the elementary identities and conventions that underlie that comparison. \hyperref[v1a-ch00-s1]{Section 0.1} treats matrix traces. \hyperref[v1a-ch00-s2]{Section 0.2} passes from orthonormal eigenfunctions to an integral kernel, derives the diagonal trace identity, and extends it to a general matrix. \hyperref[v1a-ch00-s5]{Section 0.3} distinguishes the trace of a kernel from a relative pairing and a weighted diagonal integral.
